\section{System Architecture}\label{sec:architecture}

\subsection{Hardware platform}\label{subsec:platform}

The target is the Xilinx Alveo U55C (XCVU55C-FSVH2892-2L-E): 1,303,680 LUT6, 2,607,360 flip-flops, 9,024 DSP slices, 4,032 BRAM\textsubscript{18K} blocks, 960 UltraRAM blocks, and 16\,GB of HBM2 across 32 pseudo-channels~\cite{xilinx2023vitis}. For the decoder kernels described here, the relevant resources are the logic fabric and block RAM; the HBM2 subsystem, central to a companion statevector-simulator kernel referenced but out of scope for this paper, is not used.

\subsection{Interface: AXI-Lite control, one AXI-master result port}\label{subsec:axilite}

Each kernel accepts a packed error vector over a single AXI-Lite slave port. \textbf{This revision's interface is not AXI-Lite-only}, unlike an earlier version of this design: reading a kernel's result through the AXI-Lite \texttt{ap\_return} register turns out to require either raw BAR4 PCIe-BAR memory access (blocked by filesystem permissions in every environment this was tested in, and root-owned by design on at least one host regardless of group membership) or an XRT buffer-object path that fails outright against a scalar return register, because there is no AXI4 master port for a buffer object to attach to. The fix, and the one this revision uses, is a single-element AXI master (\texttt{m\_axi}) output argument: the kernel writes its packed result to a one-word buffer that a host can read back through a standard, privilege-independent XRT buffer object. This is a small, targeted addition (one word, one AXI master interface), not a return to the HBM-pipeline architecture rejected below, but it is a real architectural change from the fully AXI-Lite design, and Section~\ref{sec:results} reports its measured HLS-level cost directly rather than describing the kernel as interface-free.

\subsection{Monolithic versus multi-kernel pipeline}\label{subsec:arch-choice}

Two standard architectures exist for a multi-stage FPGA QEC pipeline: separate compute units for encoding, syndrome extraction, and decoding, communicating through HBM ping-pong buffers; or one monolithic kernel holding all state in registers and on-chip memory. For a distance-3 code, where the entire correction state is under 20 bits, routing that state through HBM (with a per-hop round trip on the order of 100\,ns) would dominate the decode latency by roughly an order of magnitude while consuming HBM pseudo-channels a companion simulator kernel would otherwise use. The monolithic design keeps all state in registers and BRAM and is the architecture used throughout this paper. For distance-5 and larger codes, where the syndrome LUT no longer fits in on-chip memory (Table~\ref{tab:lut-scaling}) and iterative decoders require persistent message state, this calculus changes and a separate decoder kernel becomes the better design; that regime is out of scope for the present distance-3 study.
